% Complete Introduction field recovered from incomplete structured response. Not a finished source-checked package.
\section{Introduction}
\label{sec:intro}

The ion temperature \Ti{} at the edge of an H-mode plasma changes on several time scales. On KSTAR, charge-exchange spectroscopy (CES) has shown the \Ti{} pedestal building up over tens to hundreds of milliseconds after the L--H transition and falling during ELM suppression by resonant magnetic perturbations~\citep{Ko2015IonTemperature}. On ASDEX Upgrade, a dedicated sub-millisecond system has recorded the edge impurity temperature dropping and recovering within a few milliseconds of type-I ELM crashes~\citep{Viezzer2018IonHeat}; such intra-ELM dynamics lie below the 10\,ms grid used here. Requirement studies for burning-plasma devices ask for time resolution in a similar range: specifications developed for FIRE call for edge \Ti{} profiles at 10\,ms and about 10\,\% precision~\citep{Young2008PlasmaMeasurements}, and the ITER specifications list edge \Ti{} at 100\,ms and 10\,\%~\citep{Donne2007ChapterDiagnostics}. These are design targets for other devices, not demonstrated performance and not KSTAR requirements. Charge-exchange \Ti{} has also served directly in feedback: on JT-60U a core \Ti{} difference between two chords was regulated with neutral-beam power inside a 10\,ms control cycle~\citep{Yoshida2009RealTime}. Monitoring on such time scales is helped by an estimate that exists when a frame is taken. This paper examines one such capability: a causal \emph{nowcast} of the reported carbon-impurity CES temperature on one pedestal-top chord of the KSTAR 2022 campaign, selected by the data provider.

CES infers the temperature and rotation of an impurity ion from the Doppler width and shift of a line excited by charge exchange with neutral-beam atoms~\citep{Fonck1984DeterminationPlasma}; the KSTAR instrument observes the C\,VI line near 529\,nm with a separate background view~\citep{Ko2010ChargeExchange}. In the beam-modulation analysis described for KSTAR, the background of a beam-on frame is interpolated from the beam-off frames that bracket it~\citep{Lee2020TwoGaussian,Lee2026NeuralNetworkbased}, and the authors of the former argue that such an analysis cannot start before the following beam-off frame. That is an argument about the procedure, not a measurement: we did not record when the reports in our files became available, and we assign them no delay. The reports are, moreover, conventional labels rather than independent physical truth. Charge-exchange spectra superpose active and passive emission that the analysis must separate~\citep{Viezzer2011PassiveEdge}. Main-ion and carbon temperatures have been measured to differ in the edge of DIII-D, by tens of eV inside the pedestal top in H-mode and by more near the separatrix~\citep{Haskey2018MainIon}. Occasional fit failures have been associated with abrupt outliers in KSTAR CES profiles~\citep{Kim2024KineticProfile}. On the chord studied here some frames carry no usable report, and the reports that exist scatter from frame to frame (\S\ref{sec:data}). The simplest nowcast holds the last report.

One established response works on the spectrum itself. Neural networks trained on conventionally fitted spectra have long been used to read charge-exchange line parameters. At JET~\citep{Bishop1993AutomaticAnalysis,Svensson1999AnalysisJet} a monitoring demonstration delivered \Ti{} profiles at the 50\,ms exposure cadence, and a recent example comes from HL-3~\citep{Tian2026RealTimeCxrs}. At DIII-D a simplified real-time fit is reported to deliver values every 5\,ms with a delay of a few milliseconds; its agreement with the offline automatic analysis is lower on edge chords than on core chords~\citep{Chrystal2026RealTime}. On KSTAR, a two-Gaussian fit~\citep{Lee2020TwoGaussian} and the NN-CES network~\citep{Lee2026NeuralNetworkbased} both dispense with beam-off frames. NN-CES is a per-chord network trained on modulation-analysis results that maps a single beam-on spectrum to Doppler shift and width. Its authors report close agreement with the modulation analysis inside $R\approx2.2$\,m and weaker agreement farther out, where they point to strong passive emission, and they name real-time CES as a future goal. We have not examined how NN-CES behaves on the chord studied here. All of these methods read a spectrum that has been recorded, and their reference is the conventional fit, not an independent measurement.

Spectrum reading and temporal nowcasting can therefore be viewed as two stages that may be coupled. An upstream reader such as NN-CES turns a measured spectrum into a report~\citep{Lee2026NeuralNetworkbased}; a downstream estimator uses the history of such reports. Ours takes strictly past CES reports together with current and past BES and ECEI snapshots, and estimates the current report before that report is available; it does not need the current CES spectrum. This is nowcasting of a pending report. It is not forecasting of a future physical temperature or profile, as in profile prediction about 200\,ms ahead from present profiles, plasma history and proposed actuator trajectories~\citep{Abbate2021DatadrivenProfile}. A faster upstream reader would refresh the anchors on which the nowcast rests, and could equally reduce the marginal need for it; which effect prevails has to be measured. We did not implement this coupling: the evaluation uses conventional CES reports as supplied, and the connection to NN-CES is an untested proposal.

A second body of work estimates slow or missing measurements from other signals. A real-time kinetic-profile reconstruction at DIII-D takes charge-exchange data among its inputs and is trained with input dropout, so that it still returns profiles, including an ion temperature profile, when those data are absent; its ion temperature in that regime was assessed qualitatively~\citep{Shousha2023MachineLearningbased}. Faster diagnostics have been mapped to Thomson-scattering electron density and temperature at DIII-D~\citep{Jalalvand2025MultimodalSuperresolution} and COMPASS~\citep{Imrek2026TemporalSuperresolution} without the slow diagnostic as an input. Agreement was checked at the available Thomson samples, so the behaviour between samples rests on consistency arguments rather than on independent high-rate measurements. On TCV, a Kalman observer built on a transport model merges fast interferometry with slower Thomson scattering, using the slow diagnostic to correct the fast one between samples~\citep{Pastore2023ModelBased}. We do not suggest that this literature is noncausal or restricted to electron quantities. The present study asks a differently timed question about one measured ion-temperature chord, with past reports of the same diagnostic as anchors. Its fast inputs are BES~\citep{Lampert2015CombinedHydrogen} and ECEI~\citep{Yun2014Quasi3d} voltage snapshots on the 10\,ms grid. BES light depends mainly on density, and ECEI is used chiefly for relative electron-temperature fluctuations, but modelling for ASDEX Upgrade indicates that edge ECEI can also respond to density~\citep{Vanovac2018EffectsDensity}. We therefore use the snapshots as uncalibrated predictors, not as calibrated electron temperature or density and not as full-rate turbulence measurements.

The paper is organized around testable comparisons:
\begin{enumerate}\setlength{\itemsep}{1pt}\setlength{\parskip}{0pt}
\item Which causal estimators are useful when scored against hold, a causal Gaussian-process estimator and two-sided interpolation of the monotone piecewise-cubic family~\citep{fritsch1980}, which needs the next report and so serves as an offline reference (\S\ref{sec:res:best})?
\item What information do BES and ECEI add to the CES history (\S\ref{sec:res:fast})?
\item Do the answers persist under four validation designs, which are given equal weight, and does the computation fit within a 10\,ms frame (\S\ref{sec:res:designs})?
\end{enumerate}
Optional application analyses (\S\ref{sec:res:use}) and the conditions under which the estimator does not help (\S\ref{sec:res:limits}) follow. Two targets are reported. F is the CES report of the current frame. L5 is a level computed offline as the median of at least three of the five subsequent reports. Both are targets only: neither the current F nor the future-derived L5 is an input available online.

The estimators are assembled from familiar models, and no algorithmic priority is claimed; the contribution is an application-specific empirical evaluation. The results concern one chord and one reused campaign, and no untouched data exist, so the validation intervals call for corresponding caution. Meeting a computation budget is not a real-time deployment, and no physical mechanism is identified.
